\section{Метод решения}
Программа реализует межпроцессное взаимодействие по схеме «родитель-потомок» с использованием двух неименованных каналов (pipe). Архитектура программы следует паттерну абстракции операционной системы, где системно-зависимые вызовы инкапсулированы в отдельный модуль.

{\bfseries Алгоритм работы:}
\begin{enumerate}
    \item Родительский процесс запрашивает у пользователя имя выходного файла
    \item Создаются два канала: pipe1 (родитель → потомок) и pipe2 (потомок → родитель)
    \item С помощью системного вызова \texttt{fork()} создается дочерний процесс
    \item В дочернем процессе:
    \begin{itemize}
        \item stdin перенаправляется на чтение из pipe1
        \item stdout перенаправляется в выходной файл
        \item stderr перенаправляется на запись в pipe2
        \item Загружается и выполняется программа child
    \end{itemize}
    \item Родительский процесс читает строки от пользователя и отправляет их в pipe1
    \item Дочерний процесс проверяет каждую строку:
    \begin{itemize}
        \item Если строка оканчивается на '.' или ';' → выводит в stdout (записывается в файл)
        \item Иначе → выводит сообщение об ошибке в stderr (отправляется в pipe2)
    \end{itemize}
    \item Родительский процесс читает сообщения об ошибках из pipe2 и выводит их пользователю
    \item После завершения ввода (Ctrl+D) родитель ожидает завершения дочернего процесса
\end{enumerate}

\section{Описание программы}
Программа разделена на несколько модулей для обеспечения модульности и абстракции платформы.

{\bfseries Файловая структура:}
\begin{itemize}
    \item \texttt{parent.c} — основной процесс, управляет созданием потомка и каналов
    \item \texttt{child.c} — дочерний процесс, выполняет проверку строк
    \item \texttt{os.h} — абстрактный интерфейс системных вызовов
    \item \texttt{os\_linux.c} — конкретная реализация для Linux
\end{itemize}

{\bfseries Основные типы данных:}
\begin{verbatim}
typedef int pipe_t;  // Тип для дескриптора канала
typedef int pid_t;   // Тип для идентификатора процесса
\end{verbatim}

{\bfseries Функции интерфейса os.h:}
\begin{itemize}
    \item \texttt{CreatePipe()} — создание неименованного канала
    \item \texttt{CloneProcess()} — создание дочернего процесса
    \item \texttt{Exec()} — замена образа процесса
    \item \texttt{LinkFDtoIN()/LinkFDtoOUT()/LinkFDtoERR()} — перенаправление потоков
    \item \texttt{OpenObject()} — открытие файла
    \item \texttt{WritePipe()/ReadPipe()} — запись/чтение из канала
    \item \texttt{WaitProcess()} — ожидание завершения процесса
    \item \texttt{ClosePipe()} — закрытие дескриптора
\end{itemize}

{\bfseries Используемые системные вызовы Linux:}
\begin{itemize}
    \item \texttt{pipe()} — создание канала для межпроцессного взаимодействия
    \item \texttt{fork()} — создание идентичной копии процесса
    \item \texttt{dup2()} — дублирование файловых дескрипторов для перенаправления потоков
    \item \texttt{execv()} — замена текущего образа процесса новым
    \item \texttt{open()} — открытие/создание файла
    \item \texttt{read()/write()} — чтение/запись данных
    \item \texttt{close()} — закрытие файлового дескриптора
    \item \texttt{wait()} — ожидание изменения состояния процесса
\end{itemize}

{\bfseries Обработка ошибок:}
Программа проверяет возвращаемые значения всех системных вызовов и функций. При обнаружении ошибок выводится понятное сообщение с помощью \texttt{perror()} и программа завершается с ненулевым кодом возврата.